\documentclass[11pt,a4paper]{article}

\usepackage[utf8]{inputenc}
\usepackage[T1]{fontenc}
\usepackage[spanish,es-tabla]{babel}
\usepackage[margin=2.2cm]{geometry}

% Tipografía relajada, moderna y altamente legible con peso unificado
\usepackage{charter}                     % Bitstream Charter: fuente abierta, descansada y elegante (peso regular y bold neutro)
\usepackage[scaled=0.92]{beramono}       % Bera Mono para código claro y bien proporcionado

% Microtipografía y control de espaciado interlineal (leading más abierto y descansado)
\usepackage{microtype}
\usepackage{setspace}
\setstretch{1.18}                        % Leading holgado para lectura cómoda sin cansar la vista

\usepackage{booktabs}
\usepackage{tabularx}
\usepackage{enumitem}
\usepackage{titlesec}
\usepackage{listings}
\usepackage{xcolor}
\usepackage[colorlinks=true,linkcolor=blue!75!black,urlcolor=blue!75!black]{hyperref}

% Paleta de colores sobria, suave y descansada
\definecolor{primary}{RGB}{24, 48, 77}
\definecolor{secondary}{RGB}{42, 98, 148}
\definecolor{danger}{RGB}{165, 38, 38}
\definecolor{codebg}{RGB}{248, 250, 252}
\definecolor{bordercolor}{RGB}{218, 224, 233}

% Estilo de página limpio sin cabeceras ni pie de página
\pagestyle{plain}

% Estilo de títulos con la misma familia tipográfica (Charter) y peso bold homogéneo
\titleformat{\section}{\large\bfseries\color{primary}}{\thesection.}{0.5em}{}[\titlerule]
\titleformat{\subsection}{\normalsize\bfseries\color{secondary}}{\thesubsection.}{0.4em}{}
\titlespacing*{\section}{0pt}{14pt}{6pt}
\titlespacing*{\subsection}{0pt}{10pt}{5pt}

% Espaciado de listas con mayor aire
\setlist{itemsep=3.5pt, topsep=3.5pt, parsep=1pt}

% Configuración de código con tipografía clara y leading neutro
\lstset{
    basicstyle=\ttfamily\small\setstretch{1.02},
    backgroundcolor=\color{codebg},
    frame=single,
    rulecolor=\color{bordercolor},
    breaklines=true,
    breakatwhitespace=true,
    showstringspaces=false,
    keywordstyle=\color{secondary}\bfseries,
    commentstyle=\color{gray}\itshape,
    stringstyle=\color{danger},
    aboveskip=8pt,
    belowskip=8pt
}

\begin{document}

% Encabezado
\noindent
{\LARGE\textbf{\color{primary}Reporte de Seguridad: Hosting de Medios en DMOJ}}\\[0.25em]
{\large\textbf{\color{secondary}Análisis de Vulnerabilidades, Embedding Seguro de PDF y Plan de Acción}}\\[0.4em]
{\small\color{gray}Octubre 2026}
\vspace{0.4em}
\hrule height 0.8pt
\vspace{0.9em}

\section{Análisis de Amenazas y Vulnerabilidades por Formato}

El soporte de archivos multimedia en DMOJ introduce vectores de ataque diferenciados según la estructura interna de cada formato:

\subsection{Formato PNG (Portable Network Graphics)}
Formato gráfico basado en bloques (\textit{chunks}) binarios (\texttt{IHDR}, \texttt{IDAT}, \texttt{IEND}) y bloques auxiliares de metadatos (\texttt{tEXt}, \texttt{zTXt}, \texttt{iTXt}).
\begin{itemize}
    \item \textbf{Decompression Bombs (Pixel Floods):} Un atacante manipula la cabecera \texttt{IHDR} declarando dimensiones gigantescas (e.g. $50000 \times 50000$ px) en un archivo comprimido de pocos kilobytes. Al procesarlo, la librería del servidor o el navegador del usuario asigna memoria para cada píxel sin comprimir (RGBA $\approx$ 10 GB de RAM), desatando un colapso por \textit{Out-Of-Memory} (DoS).
    \item \textbf{MIME-Sniffing y XSS Almacenado:} Los chunks de metadatos permiten inyectar texto arbitrario. Si el servidor responde sin \texttt{X-Content-Type-Options: nosniff}, navegadores antiguos o visores heurísticos detectan cadenas HTML/JS dentro del PNG y las ejecutan en el contexto del sitio.
    \item \textbf{Ataques Polyglot (PNG + HTML):} Archivos con cabecera válida de PNG que contienen etiquetas HTML válidas capaces de ser interpretadas por navegadores si se invocan directamente.
\end{itemize}

\subsection{Formato JPG / JPEG (Joint Photographic Experts Group)}
Formato ráster que utiliza segmentos de aplicación (\texttt{APP0} a \texttt{APP15}) para almacenar metadatos en estándares EXIF, IPTC y XMP, delimitado por el marcador de inicio \texttt{0xFFD8FF} y cierre \texttt{0xFFD9} (\texttt{EOI}).
\begin{itemize}
    \item \textbf{Inyección en Metadatos EXIF:} Atacantes inyectan código malicioso (e.g. \texttt{"\textless script\textgreater alert(1)\textless/script\textgreater"}) en campos estándar como \texttt{Artist}, \texttt{CameraModel} o \texttt{UserComment}. Si una vista de administración o foro renderiza estos valores sin escapar, ocurre un XSS almacenado.
    \item \textbf{Progressive JPEG DoS:} Imágenes con múltiples pasadas espectrales entrelazadas diseñadas para forzar ciclos de cálculo excesivos en los decodificadores gráficos.
    \item \textbf{Cargas útiles post-EOI:} Los visualizadores ignoran datos situados después del marcador \texttt{0xFFD9}, permitiendo anexar scripts o binarios ocultos que evaden escaneos superficiales.
\end{itemize}

\subsection{Formato PDF (Portable Document Format)}
Documento estructurado en objetos con soporte nativo para flujos de datos comprimidos, acciones automáticas y motores de scripting internos. \textbf{Es el formato de mayor riesgo para los usuarios finales.}
\begin{itemize}
    \item \textbf{Ejecución de JavaScript Activo (\texttt{/JS}, \texttt{/JavaScript}):} Los visores integrados en navegadores modernos (PDFium en Chromium, PDF.js en Firefox) soportan JavaScript. Si un PDF se despacha de forma \textit{inline} en el mismo dominio (\texttt{judge.com/media/archivo.pdf}), el script corre en el origen del sitio y puede \textbf{robar cookies de sesión (\texttt{sessionid}) y tokens CSRF}.
    \item \textbf{Acciones de Lanzamiento y Formularios (\texttt{/Launch}, \texttt{/SubmitForm}):} Permiten intentar abrir procesos en la máquina del usuario o enviar datos locales a servidores externos de control.
    \item \textbf{Archivos Incrustados Ocultos (\texttt{/EmbeddedFiles}):} Embebido de malware o ejecutables destinados a engañar a los concursantes que descargan el material.
\end{itemize}

\subsection{Formato ZIP (PKZIP Archive)}
Estructura de archivo contenedor con cabeceras locales (\texttt{PK\textbackslash x03\textbackslash x04}) y un directorio central. Utilizado en DMOJ tanto para paquetes de evaluación como para archivos de plantilla/starter kits que los concursantes descargan en sus máquinas.
\begin{itemize}
    \item \textbf{RCE en el Servidor vía Zip Slip:} Rutas con secuencias relativas (e.g. \texttt{../../dmoj/settings.py}). Al descomprimirse en el backend, sobrescribe archivos de la aplicación o del sistema operativo, logrando ejecución remota de código en el servidor.
    \item \textbf{RCE en el Cliente vía Plantillas Maliciosas:} Cuando un creador de problemas o atacante sube un ZIP con código de plantilla para los participantes, puede incluir archivos que logran ejecución de código en la máquina del concursante al descargarlo o abrirlo:
    \begin{itemize}
        \item Binarios o scripts directos (\texttt{.exe}, \texttt{.sh}, \texttt{.bat}, \texttt{.vbs}, \texttt{.ps1}).
        \item Archivos de configuración de IDEs (e.g. \texttt{.vscode/tasks.json}, subcarpetas \texttt{.idea}) que ejecutan comandos automáticos al abrir el proyecto en VS Code o CLion.
        \item Scripts de build maliciosos (\texttt{Makefile}, \texttt{CMakeLists.txt}) o archivos con extensiones dobles engañosas (\texttt{template.cpp.exe}).
    \end{itemize}
    \item \textbf{Zip Bomb (Compresión Recursiva / DoS Extremo):} Ratios de compresión anómalos (ej. 50 KB expanden a 50 GB) que saturan el disco o la RAM al descomprimirse.
    \item \textbf{Manipulación de Symlinks:} Enlaces simbólicos apuntando a rutas del sistema operativo (\texttt{/etc/passwd}, \texttt{db.sqlite3}) para exfiltrar datos del servidor o del equipo de la víctima.
\end{itemize}

\vspace{0.4em}
\begin{table}[h!]
\centering
\small
\setstretch{1.08}
\caption{Matriz Comparativa de Riesgos y Contramedidas}
\label{tab:riesgos}
\begin{tabularx}{\textwidth}{l p{4.2cm} p{3.8cm} X}
\toprule
\textbf{Formato} & \textbf{Vectores Críticos} & \textbf{Impacto Potencial} & \textbf{Contramedida Técnica Obligatoria} \\
\midrule
\textbf{PNG} & Decompression bomb, metadatos con scripts, polyglot & Consumo de RAM (DoS), XSS por sniffing & Re-codificación con Pillow, límite \texttt{MAX\_IMAGE\_PIXELS}, header \texttt{nosniff}. \\
\addlinespace
\textbf{JPG} & Payloads en etiquetas EXIF, datos post-EOI & XSS en vistas administrativas & Eliminación total de metadatos EXIF, re-muestreo a búfer nuevo en memoria. \\
\addlinespace
\textbf{PDF} & Scripts activos (\texttt{/JS}), acciones (\texttt{/Launch}), iframes & Robo de sesión (\texttt{sessionid}), phishing & Aislamiento de origen, \texttt{\textless iframe\textgreater} restringido o rasterización server-side. \\
\addlinespace
\textbf{ZIP} & Zip Slip (\texttt{../}), RCE en cliente (scripts/templates), symlinks & RCE en servidor o en estación del usuario, DoS & Lista blanca estricta de extensiones de código, normpath, ratio máx 20:1, sin symlinks. \\
\bottomrule
\end{tabularx}
\end{table}

\section{Análisis Especializado: Previsualización y Embedding Seguro de PDFs}

Forzar la descarga exclusiva (\texttt{Content-Disposition: attachment}) es la medida más restrictiva contra ataques a clientes, pero destruye la experiencia de usuario (UX) cuando se requiere previsualizar enunciados de problemas en la web.

El debate técnico sobre cómo incrustar PDFs no confiables (e.g., discusión técnica de referencia en \textit{Security StackExchange} / PrivateBin / MDN) expone trade-offs determinantes:

\subsection{Comparativa de Mecanismos Frontend: \texttt{\textless embed\textgreater}, \texttt{\textless object\textgreater} e \texttt{\textless iframe\textgreater}}
\begin{enumerate}[leftmargin=*]
    \item \textbf{Elemento \texttt{\textless embed\textgreater} (Inseguro - Descartado):} Históricamente utilizado para incrustar plugins. En varios navegadores y contextos de renderizado, el contenido de \texttt{\textless embed\textgreater} puede interactuar o heredar acceso directo al documento padre (\texttt{top-level document}), carece de control granular de sandbox y no admite directivas modernas de carga (\texttt{loading="lazy"}, \texttt{referrerpolicy}).
    \item \textbf{Elemento \texttt{\textless object\textgreater} (No recomendado):} Requiere relajar la directiva CSP \texttt{object-src}. Las mejores prácticas de seguridad web exigen fijar \texttt{object-src 'none'} para neutralizar la ejecución de plugins Flash/Java/ActiveX heredados. Usar \texttt{\textless object\textgreater} debilita el perfil global de seguridad de la aplicación.
    \item \textbf{Elemento \texttt{\textless iframe\textgreater} (Estándar Recomendado):}
    \begin{itemize}
        \item Se gobierna mediante la directiva CSP \texttt{frame-src}, permitiendo autorizar exclusivamente el subdominio aislado de medios.
        \item Admite aislamiento de contexto y carga perezosa.
        \item \textbf{El dilema del atributo \texttt{sandbox}:} Aplicar un sandbox ciego (e.g. \texttt{sandbox=""} sin \texttt{allow-scripts}) suele romper el visor de PDF integrado en el navegador (como Mozilla PDF.js o Chrome PDFium), ya que estos motores requieren JavaScript interno para renderizar el documento. No obstante, los visores nativos modernos ya cuentan con su propio sandboxing de procesos a nivel de sistema operativo.
        \item \textbf{Mitigación óptima en iframe:} La protección crítica consiste en que el \texttt{src} del iframe apunte a un \textbf{origen cruzado sin cookies} (\texttt{https://media.dmoj.local/pdf/...}), configurando \texttt{sandbox="allow-scripts allow-same-origin"} pero \textbf{omitiendo} \texttt{allow-top-navigation} y \texttt{allow-popups} para impedir secuestro de ventana o ataques de phishing.
    \end{itemize}
\end{enumerate}

\subsection{La Solución de Máxima Seguridad: Rasterización Server-Side (PDF-to-Image)}
Para eliminar por completo el riesgo de explotación en el cliente (vulnerabilidades 0-day en PDFium/PDF.js o bypasses de sandbox), la estrategia más robusta para DMOJ consiste en renderizar una previsualización en el servidor:
\begin{itemize}
    \item Un worker asíncrono o servicio aislado en backend (utilizando \texttt{pdftoppm} de Poppler o MuPDF encapsulado en un contenedor mínimo con límites de memoria y CPU) convierte cada página del PDF en una imagen ráster estática (\texttt{.png}).
    \item La interfaz web de DMOJ muestra la previsualización del enunciado mediante etiquetas inertes \texttt{\textless img\textgreater}.
    \item Esto reduce la superficie de ataque del cliente a \textbf{cero} (no hay ejecución de JavaScript de PDF, no hay manipulación de DOM ni exploits de visores), reservando el archivo PDF original como enlace de descarga opcional bajo \texttt{Content-Disposition: attachment}.
\end{itemize}

\section{Arquitectura Defensiva Global (Defense-in-Depth)}

Se articulan tres capas de defensa simultáneas para blindar la plataforma:

\subsection{1. Aislamiento de Origen (Dominio sin Cookies)}
El vector más destructivo en el cliente (robo de cookies y XSS) queda neutralizado aislando el origen:
\begin{itemize}
    \item El contenido de medios debe servirse desde un subdominio dedicado (ej. \texttt{media.dmoj.local}).
    \item En la configuración de Django (\texttt{settings.py}), se delimita el alcance de las cookies de sesión al dominio principal:
    \begin{lstlisting}[language=Python]
SESSION_COOKIE_DOMAIN = "dmoj.local"  # No se transmite al dominio de medios
SESSION_COOKIE_HTTPONLY = True
SESSION_COOKIE_SECURE = True
    \end{lstlisting}
    \item Si un archivo lograra ejecutar JavaScript en su propio visor, correrá en un origen sin credenciales ni acceso al DOM del juez.
\end{itemize}

\subsection{2. Encabezados HTTP de Seguridad Estricta}
El servidor web (Nginx) y las respuestas de Django deben emitir encabezados restrictivos:
\begin{itemize}
    \item \textbf{X-Content-Type-Options (\texttt{nosniff}):} Impide que el navegador ignore el MIME declarado e interprete HTML oculto en imágenes o documentos.
    \item \textbf{Content-Security-Policy:}
    \begin{itemize}
        \item En el sitio principal: \texttt{object-src 'none'; frame-src https://media.dmoj.local;}
        \item En el dominio de medios: \texttt{default-src 'none'; style-src 'unsafe-inline'; script-src 'none'; sandbox;}
    \end{itemize}
    \item \textbf{Content-Disposition:}
    \begin{itemize}
        \item Imágenes (\texttt{.png}, \texttt{.jpg}): \texttt{inline} (para visualización en enunciados).
        \item ZIPs: \texttt{attachment; filename="seguro.zip"} (descarga obligatoria).
        \item PDFs: \texttt{inline} únicamente dentro de \texttt{\textless iframe\textgreater} en el dominio de medios aislado; en cualquier otro contexto, forzar \texttt{attachment}.
    \end{itemize}
    \item \textbf{Cross-Origin-Resource-Policy (\texttt{cross-origin}):} Permite cargar imágenes en los enunciados mediante etiquetas \texttt{\textless img\textgreater} sin comprometer la política de aislamiento.
\end{itemize}

\subsection{3. Almacenamiento Seguro, Sanitización y Control de Templates}
\begin{itemize}
    \item \textbf{Identificadores Opacos (UUIDv4):} Los ficheros se guardan en el disco con identificadores aleatorios: \texttt{problem\_media/\textless problem\_code\textgreater/\textless uuid4\textgreater.\textless ext\textgreater}. El nombre legible del usuario solo se utiliza en la cabecera \texttt{Content-Disposition} debidamente escapado.
    \item \textbf{Filtrado Estricto de Contenido de ZIPs (Anti-RCE Cliente y Servidor):} Todo paquete ZIP (ya sea de datos o plantillas para participantes) debe someterse a inspección profunda de sus entradas:
    \begin{itemize}
        \item \textbf{Lista blanca de extensiones:} Permitir únicamente archivos de código fuente y datos puros (\texttt{.cpp}, \texttt{.c}, \texttt{.h}, \texttt{.hpp}, \texttt{.py}, \texttt{.java}, \texttt{.in}, \texttt{.out}, \texttt{.txt}).
        \item \textbf{Bloqueo de directorios ocultos y de configuración:} Rechazar cualquier entrada que comience por punto o rutas como \texttt{.vscode/}, \texttt{.idea/}, \texttt{.git/}, impidiendo ejecución arbitraria al abrir el proyecto en un IDE.
        \item \textbf{Prohibición total de ejecutables y scripts:} Rechazar automáticamente extensiones de scripts o binarios (\texttt{.sh}, \texttt{.bat}, \texttt{.ps1}, \texttt{.exe}, \texttt{.dll}, etc.) y scripts de build no confiables (\texttt{Makefile}).
    \end{itemize}
    \item \textbf{Punto de Montaje No Ejecutable:} El directorio de medios en el servidor debe montarse con flags \texttt{noexec, nosuid, nodev} en \texttt{/etc/fstab}, impidiendo la ejecución de cualquier script que lograse filtrarse.
    \item \textbf{Re-codificación en Memoria:} Las imágenes son deserializadas y re-exportadas limpias con Pillow, destruyendo cualquier payload alojado en metadatos o chunks no esenciales antes de tocar el disco.
\end{itemize}

\section{Plan de Acción e Implementación}

Hoja de ruta secuencial para habilitar el hosting de medios en DMOJ:

\begin{enumerate}[leftmargin=*]
    \item \textbf{Fase 1: Corrección Inmediata de Encabezados y Despacho (Semana 1)}
    \begin{itemize}
        \item Modificar \texttt{problem\_media\_file} en \texttt{judge/views/problem\_data.py}:
        \begin{itemize}
            \item Reemplazar \texttt{application/octet-stream} indiscriminado por el MIME exacto.
            \item Forzar cabeceras \texttt{X-Content-Type-Options: nosniff} y \texttt{Content-Security-Policy}.
            \item Asignar \texttt{Content-Disposition: attachment} para archivos \texttt{.zip}.
        \end{itemize}
        \item Restringir el acceso a medios verificando permisos de visualización del problema (\texttt{is\_editable\_by} o problema publicado no confidencial).
    \end{itemize}

    \item \textbf{Fase 2: Pipeline de Validación y Desinfección (Semanas 2 y 3)}
    \begin{itemize}
        \item Integrar el validador provisto (\texttt{security/media\_security\_validator.py}) en el método \texttt{clean\_upload} de \texttt{ProblemMediaUploadForm}.
        \item Activar el re-muestreo obligatorio de imágenes PNG/JPG descartando metadatos EXIF.
        \item Habilitar la validación preventiva de archivos ZIP: verificación de rutas seguras (\texttt{os.path.normpath} sin \texttt{..}), lista blanca estricta de extensiones de código/texto para plantillas, bloqueo de carpetas ocultas de IDEs (\texttt{.vscode}, \texttt{.idea}), rechazo total de symlinks y límite de ratio a 20:1.
        \item Bloquear subida de PDFs que contengan patrones \texttt{/JavaScript}, \texttt{/JS} o \texttt{/Launch}.
    \end{itemize}

    \item \textbf{Fase 3: Aislamiento de Infraestructura y Preview de PDF (Semanas 4 y 5)}
    \begin{itemize}
        \item Desplegar la configuración de Nginx (\texttt{security/nginx\_media\_hardening.conf}) usando \texttt{X-Accel-Redirect} en el subdominio dedicado \texttt{media.dmoj.local}.
        \item Implementar el preview de PDFs:
        \begin{itemize}
            \item \textit{Estrategia recomendada:} Generación de miniaturas PNG de las páginas del PDF mediante worker aislado con \texttt{pdftoppm}.
            \item \textit{Estrategia directa en visor:} Incrustar exclusivamente mediante \texttt{\textless iframe\textgreater} apuntando a \texttt{media.dmoj.local}, con CSP \texttt{frame-src} estricto y descartando terminantemente \texttt{\textless embed\textgreater} u \texttt{\textless object\textgreater}.
        \end{itemize}
        \item Montar directorio físico con permisos restringidos (\texttt{chmod 750}) y flag \texttt{noexec}.
    \end{itemize}

    \item \textbf{Fase 4: Testing Automatizado y Despliegue (Semana 6)}
    \begin{itemize}
        \item Integrar la suite de pruebas unitarias (\texttt{security/test\_security\_pipeline.py}) al pipeline de CI/CD para garantizar que futuras modificaciones no introduzcan regresiones de seguridad.
        \item Habilitar tokens temporales de descarga firmados con HMAC para problemas en concursos activos con enunciados privados.
    \end{itemize}
\end{enumerate}

\section{Configuración Esencial de Referencia (Nginx Hardening)}

Directiva central que debe incorporarse en el servidor web para el dominio de medios:

\begin{lstlisting}[language=sh]
location /internal_media/ {
    internal;
    alias /var/dmoj/problem_media/;

    # Bloqueo total de MIME-sniffing y aislamiento CSP
    add_header X-Content-Type-Options "nosniff" always;
    add_header Content-Security-Policy "default-src 'none'; style-src 'unsafe-inline'; sandbox allow-scripts allow-same-origin;" always;
    add_header X-Frame-Options "SAMEORIGIN" always;
    add_header Cross-Origin-Resource-Policy "cross-origin" always;

    # Imagenes seguras se muestran inline en enunciados
    location ~* \.(png|jpg|jpeg)$ {
        add_header Content-Disposition "inline" always;
    }
    # Archivos ZIP forzados a descarga obligatoria
    location ~* \.zip$ {
        add_header Content-Disposition "attachment" always;
    }
    # PDFs servidos para previsualizacion aislada en iframe
    location ~* \.pdf$ {
        add_header Content-Disposition "inline" always;
    }
}
\end{lstlisting}

\end{document}
